\section{Clique interfaces, block trees and capacity}\label{sec:cliques}

Fix a root $z$ of $K_n$, external contributions $A,B\ge0$ with
$A+B$ even, and a desired internal root red sum $S$. Only internal
edges are to be checked, using final root counts
$(A+S,B+2(n-1)-S)$.

\begin{theorem}[Complete clique root interface]\label{thm:cliqueport}
For every $n\ge4$ and $0\le S\le2(n-1)$ there is a legal internal
colouring with the prescribed root sum. For $n\ge5$ every nonroot
vertex may additionally have both colour degrees in $[1,2n-3]$.
For $n=4$ the same strengthening holds except at
$(A,B,S)=(1,3,4),(3,1,2)$; the ordinary interface still holds there.
\end{theorem}
\begin{proof}
For $n=4,5$, put $r=n-1$. External counts exceeding $2r$ are
comparison-equivalent to $2r+1$ for odd counts and $2r+2$ for even
counts, since every nonroot count is at most $2r$. Consequently the
entire base domain is
\[
 A,B\in\{0,\ldots,2r+2\},\quad A+B\equiv0\pmod2,
 \quad S\in\{0,\ldots,2r\}.
\]
It contains $287+549=836$ entries. Actual edge-weight witnesses are
in \code{rooted_clique_interfaces.json}. The independent standard-library
checker \code{rooted_clique_certificate.py} verifies equality with this
whole domain, root sums, every actual colour comparison, and every
nonroot bound outside the two exceptions. Thus the bases include
arbitrarily large physical external counts by exact clipping, not by
an unproved extrapolation.

For $n\ge6$ use strong induction by two vertices. If
$2\le S\le2n-4$, take a rooted $K_{n-2}$ with the same external
counts and internal sum $S-2$. Add $u,v$, set $wu=0,wv=2$ for
every old $w$, and $uv=1$. All old colour degrees increase by two;
the new counts are
\[
 u:(1,2n-3),\qquad v:(2n-3,1).
\]
Old nonroot counts are within $[2,2n-4]$, so all their new edges
are legal. Write the final root counts as $p,q$. If neither equals
$2n-3$, its two new pure edges are also legal. Otherwise both
$p,q$ are odd because $p+q$ is even. Choose an old nonroot $w_0$
and change its two new weights from $(0,2)$ to $(1,1)$. All old
counts stay fixed, while the new counts become
\[
 u:(2,2n-4),\qquad v:(2n-4,2).
\]
Inductively strengthened old nonroot counts, after adding two, lie
in $[3,2n-5]$ and are strictly separated from these boundary values.
The root counts are odd, so its comparisons with the new even values
hold. When $n=6$ and the old $K_4$ is one of the two exceptional
bases, the final root counts are instead $p=q=7$, different from
$2n-3=9$; the unmodified construction already works. Thus no missing
strengthened base is needed.

For $S=0,1$, take the old $K_{n-2}$ with the same internal sum but
formal external counts $A'=A-2,B'=B+2$. If $A\ge2$, these are
physical. If $A=0,1$, obtain the old witness using the nonnegative
replacement $\widehat A=A\bmod2$ and then reinterpret its root red
count with $A'$. This keeps the old internal edges legal: for $S=0$
the root has no red edge; for $S=1$ its sole red incidence meets a
vertex of red degree at least one, whereas the reinterpreted root
red count is nonpositive. Its blue count is unchanged. These formal
negative offsets are only an intermediate comparison argument.

Add $u,v$ with $wu=0,wv=2$ for old nonroots, $zu=zv=0$, and
$uv=2$. Relative to the reinterpreted old counts every old vertex
gains two in each colour, so the root now has exactly the originally
specified physical counts and internal sum $S$. The new counts are
$(2,2n-4),(2n-4,2)$. The old nonroot counts lie in $[3,2n-5]$;
the root blue count is at least $2n-3$. All new comparisons follow.
The old base parameter $S=0$ or 1 is never an exceptional $K_4$
parameter. Finally obtain $S=2n-3,2n-2$ by exchanging the colours
and $A,B$. These cases cover every $S$ and preserve the asserted
nonroot bounds.
\end{proof}

\subsection{Bouquets and arbitrary cores}\label{sec:cliquecapacity}

\begin{theorem}[A8]\label{thm:bouquet}
At a vertex of old degree $h$ in a given legal colouring, attach
cliques of orders $n_i\ge4$ sharing only that vertex, with total
new degree $L=\sum_i(n_i-1)\ge h$. Every old edge colour can be
preserved. Alternatively, for an arbitrary uncoloured simple core
$C$, attach such a bouquet at every vertex $v$ with
$L_v\ge d_C(v)$. The entire expanded graph can then be coloured.
\end{theorem}
\begin{proof}
For the first assertion, let the old counts at the root be $a,b$.
An added red contribution $s\in[0,2L]$ gives final counts
$(a+s,b+2L-s)$. Each old pure edge forbids at most one value of
$s$ and each mixed edge at most two. There are at most $2h$
forbidden values and $2L+1>2h$ candidates. Choose an allowed $s$
and split it into $s_i\in[0,2(n_i-1)]$. Clique $i$ receives external
counts
\[
 A_i=a+s-s_i,\quad
 B_i=b+2L-s-\bigl(2(n_i-1)-s_i\bigr).
\]
These are nonnegative with even sum. Theorem~\ref{thm:cliqueport}
colours it with internal sum $s_i$. Every clique uses the same
final root counts and all old neighbours keep their counts.

For the second assertion give every core edge weight 1. Put
$h_v=d_C(v)$ and $D_v=h_v+L_v$. In any vertex order choose
$R(v)\in[h_v,h_v+2L_v]$. A previously processed neighbour forbids
at most $R(w)$ and $2D_v-B(w)$, so at most $2h_v$ values are
forbidden. There are at least $2h_v+1$ candidates. Choose all red
targets greedily and set $B(v)=2D_v-R(v)$. Every mixed core edge
now distinguishes both colours. Split the remaining red contribution
among the cliques and realise it by the root interface theorem.
\end{proof}

The first capacity bound cannot uniformly be reduced to $L=h-1$
for $h\ge4$. Pair the integers $0,\ldots,2h-2$ within their even
and odd subsequences using adjacent values, covering a final odd
triple by two overlapping pairs. This gives exactly $h$ pairs
$(a_i,b_i)$ with $a_i=b_i+2$, distinct $a_i$, distinct $b_i$, and
union $\{0,\ldots,2h-2\}$. On a central $K_{h+1}$ with vertices
$v,u_1,\ldots,u_h$ put every weight 1. Attach a rooted $K_{h+1}$
at each $u_i$, with internal red sum $a_i$ using the theorem.
The old counts are
\[
 v:(h,h),\qquad u_i:(h+a_i,\ 3h-2-b_i).
\]
All central mixed edges are legal: both lists at the $u_i$ are
distinct and at least $h+2$. Add $K_h$ at $v$. Its contribution
$s\in[0,2h-2]$ is forbidden by $vu_i$ precisely at $s=a_i$ or
$s=b_i$. The pairs cover every candidate. Nonetheless the entire
graph is colourable by the block-tree theorem below, rooted at $v$:
each nonroot cutvertex has equal parent and child capacities. Thus
this is a sharp preserved-colour obstruction, not a graph obstruction.

\subsection{Clique block trees with full root freedom}

\begin{theorem}[A7]\label{thm:cliquetree}
Let every block of a connected graph be a clique of order at least
four. Root its block-vertex incidence tree at any vertex $z$.
At every nonroot cutvertex require the sum of its child-block degrees
$\sum(n_i-1)$ to be at least its parent-block degree $n_0-1$.
For all physical external counts $A,B$ at $z$ and all
$S\in[0,2d_G(z)]$, the graph has a legal internal colouring with
root internal red sum $S$. In particular all equal-order clique
block trees satisfy this conclusion.
\end{theorem}
\begin{proof}
First colour the bouquet of root blocks, distributing the specified
sum $S$ between them and using the root interface with the contributions
of the other root blocks included as external counts. Process the
other cutvertices outwards. Before its children are added a cutvertex
belongs to its one parent block and has current degree $h=n_0-1$.
The child bouquet has capacity at least $h$, so
Theorem~\ref{thm:bouquet} extends the current colouring without changing
any old edge weight. Actual external root counts remain included in
all old comparisons. Adding blocks away from $z$ does not change
its internal red sum. The finite incidence tree is eventually exhausted.
\end{proof}

In a vertex-minimal counterexample, any bouquet of leaf clique blocks
of order at least four at a vertex must have capacity $L<h$, where
$h$ is its degree in the remainder. Otherwise delete the leaf blocks,
colour the smaller remainder, and apply Theorem~\ref{thm:bouquet}.
If the remainder is $K_1$, use the root interface directly. If it
is $K_2$, give its edge weight 1 and the bouquet total red contribution
1; the endpoints have counts $(2,2L)$ and $(1,1)$, so the old edge
is legal and all cliques can be filled. Hence no hidden $K_2$
assumption is used in this minimality argument.

\section{Gluing coloured pieces along a tree}\label{sec:pieceglue}

Let connected pieces $P_i$ partition the edges of $G$. Let $U$ be
the vertices belonging to more than one piece; all other vertices
are called private. Suppose the bipartite incidence graph of pieces
and shared vertices is a tree. In the given local weighting of $P_i$,
write $d_i,r_i,b_i$ for the local degrees and set
\[
 a_i(v)=r_i(v)-d_i(v),\quad
 M_v=\max_{i:v\in P_i}|a_i(v)|,\quad
 I_v=[d_G(v)-M_v,d_G(v)+M_v]\quad(v\in U).
\]

\begin{theorem}[A9]\label{thm:pieceglue}
One may obtain a legal colouring solely by whole-piece colour flips
provided that: private--private edges are locally legal; each edge
from shared $v$ to private $u$ in $P_i$ satisfies
$d_G(v)-M_v>\max(r_i(u),b_i(u))$; and each shared--shared edge $uv$
satisfies $|d_G(u)-d_G(v)|>M_u+M_v$.
\end{theorem}
\begin{proof}
Root the incidence tree at a piece. On reaching a shared vertex its
parent piece's flip is already fixed. Starting with its contribution
$s=\pm a_i(v)$, choose the flip of each child piece to minimise the
new absolute sum. The identity
\[
 \min\{|s+a|,|s-a|\}\le\max\{|s|,|a|\}
\]
keeps $|s|\le M_v$. Each child piece has only one parent shared
vertex, so no choice is made twice. At the end the global red and
blue degrees at $v$ are $d_G(v)\pm s$, both in $I_v$.
Private counts are their local counts up to a swap. The three edge
types in the hypotheses are respectively preserved by local legality,
strict separation above private counts, or disjoint shared intervals.
Every edge is therefore legal.
\end{proof}

If the given pieces have maximum degree at most $q$, a sufficient
uniform condition is $d_G(v)\ge3q+1$ for shared vertices and degree
gap at least $2q+1$ on shared--shared edges. Indeed $M_v\le q$
and private colour degrees are at most $2q$. Taking known colourings
of cycle or cubic blocks gives, for example, the conditions shared
degree at least 10 and adjacent shared-degree gap at least 7. This
allows unbounded global maximum degree and pieces that are not cliques.

\subsection{Why whole-piece flips do not suffice in general}\label{sec:piecebarrier}

For each $r\ge2$, identify the distinguished left vertex $z$ of
two copies of $K_{r,r}$. In each copy put weight 2 in column 1,
weight zero in column 2, weight 1 in all remaining entries of row
$z$, and zero elsewhere. The local red sums are $r$ at $z$, 2 at
other rows, $2r$ at column 1, zero at column 2, and 1 at other
columns. As the local original degrees are equal, these values
verify local legality on every edge.

The local root counts are $(r,r)$ even after a whole-piece flip.
After identification the root counts are $(2r,2r)$, and in each
flipped or unflipped copy it still has a same-colour neighbour of
that colour degree $2r$. Thus all four choices of whole-piece flips
fail, although the incidence graph is a tree. For $r\ge3$, changing
one entry of row $z$ in a remaining column from 1 to zero in one
copy repairs the whole graph: the root becomes $(2r-1,2r+1)$,
the changed column's red sum becomes zero, and all other local
counts persist. At $r=2$ replace one copy's matrix by rows
$(0,0),(2,2)$. These direct repairs show why the obstruction concerns
given-piece gluing and not the conjecture.
